\section{Implementation in SCIP}\label{sec:implementation}

We implemented the cuts as a separator for SCIP~10.0.2 \citep{SCIP10},
written in Python on top of PySCIPOpt~6.2.1 \citep{PySCIPOpt}, with exact
arithmetic in Python fractions and SymPy, ball arithmetic in Arb through
python-flint, and HiGHS for the direction LPs. The separator adds global cuts
at the root node. It does not change SCIP's constraint handlers, presolve,
propagation or branching, and every native nonlinear constraint stays in the
model. Figure~\ref{fig:pipeline} shows the data flow. This section describes
the parts on which validity depends and the choices that shape the
computational results; Appendix~\ref{app:software} lists the work limits.

\begin{figure}[t]
\centering
\begin{tikzpicture}[
  node distance=4mm and 6mm,
  box/.style={draw,rounded corners=2pt,align=center,font=\small,inner sep=3pt,minimum height=9mm},
  trusted/.style={box,fill=gray!12},
  num/.style={box,dashed},
  >={Stealth[length=2mm]}]
\node[trusted] (src) {source model\\exact binary64};
\node[trusted,right=of src] (bld) {model builder\\DAG check, domains,\\affine bounds};
\node[box,right=of bld] (scip) {SCIP\\native constraints,\\presolve, branching};
\node[num,below=9mm of src] (disc) {block discovery};
\node[num,right=of disc] (lp) {direction LP\\over samples};
\node[trusted,right=of lp] (sup) {exact support\\for binary64\\direction};
\node[trusted,right=of sup] (elim) {elimination,\\safe export,\\stored-row check};
\draw[->] (src) -- (bld);
\draw[->] (bld) -- (scip);
\draw[->] (bld.south) -- ++(0,-3mm) -| (disc.north);
\draw[->] (disc) -- (lp);
\draw[->] (lp) -- (sup);
\draw[->] (sup) -- (elim);
\draw[->] (elim.north) |- node[pos=0.25,right,font=\small] {global cut} (scip.south east);
\draw[->,dotted] (scip.south west) ++(2mm,0) |- node[pos=0.25,left,font=\small] {LP point} (lp.north east);
\end{tikzpicture}
\caption{Data flow of the separator. Shaded steps are part of the
certificate and are repeated by replay; dashed steps only choose what to
certify and need not be correct. SCIP's own computations are outside the
certificate.}
\label{fig:pipeline}
\end{figure}


\subsection{A faithful model}\label{sec:faithful}

A cut certificate refers to the model that the solver solves. Two ordinary
steps of model loading, assembling expressions and simplifying them, can
change that model without notice. We guard against both in every mode,
including the native baseline, so that all modes solve the same model. A
third paragraph describes the bounds that certificates use.

\paragraph{Arithmetic.}
Every finite binary64 number in the source file is interpreted as its exact
rational value. Assembling a native expression can still change values: in
$(0.1x)^2$ the source coefficient is the exact square of the binary64 number
$0.1$, whereas a coefficient computed in floating point is that square
rounded again, and a constraint constructor may move a constant to the other
side with a rounded subtraction. After assembly, the builder reads back the
PySCIPOpt expression object that it passes to SCIP and compares its exact
value with the source expression by symbolic subtraction and expansion, with
no tolerance. If they differ, it rebuilds the expression as an explicit
expression graph with separate constant nodes and compares again; a remaining
difference rejects the model. Constants are kept inside the expression, so
the constructor cannot round them into the row sides, which are passed
unchanged. This is translation validation
\citep{PnueliSiegelSingerman1998,Necula2000} of the model import. It
addresses the observation of \citet[Section~20]{Neumaier2004} that modeling
systems introduce uncontrolled rounding, and the requirement of
\citet[Section~3.1]{SchichlNeumaier2005} that simplification of an expression
graph introduce no roundoff. Its boundary is the expression object;
PySCIPOpt's translation of it into SCIP's internal expression, and SCIP's
later simplification, presolve and floating-point LP, are outside it.

\paragraph{Domains.}
Simplification can silently enlarge a domain: $0\cdot\log x$, $x/x$,
$(\sqrt x)^2$ and $(\log x)^0$ are defined only for $x>0$, $x\ne0$, $x\ge0$
and $x>0$, respectively. The builder visits every node of the source
expression before simplification and proves each domain requirement from
bounds and from affine rows; for example, $x-y=1$ proves that $\log(x-y)$ is
defined, with no bounds on $x$ and $y$. Requirements that cannot be proved
are enforced by existential constraints in a new variable $u$:
\[
d\ne0\iff\exists u:\ du=1,\qquad d>0\iff\exists u\ge0:\ du=1,\qquad d\ge0\ \text{is imposed directly.}
\]
The first is the device of \citet{Rabinowitsch1930}. Over the reals these
constraints describe strict domains exactly, without a positive tolerance
such as the one used in common benchmark setups \citep{BestuzhevaEtAl2025};
SCIP then enforces them, like every constraint, with its feasibility
tolerance. A variable power $b(x)^{e(x)}$ is admitted as $\exp(e(x)\log b(x))$
only when $b(x)>0$ is proved \citep[compare][]{BelottiEtAl2009}; otherwise
the model is rejected, because imposing $b>0$ could remove points at which
the power is defined for other reasons.

\paragraph{Bounds.}
Support domains need finite bounds. Original affine rows imply bounds that
may be missing from the declared box: for a row $\sum_ia_ix_i\le b$ with
$a_j\ne0$, $a_jx_j\le b-\sum_{i\ne j}\min\{a_i\ell_i,a_iu_i\}$. Any finite
sequence of such steps, with integer rounding for integer variables, yields
bounds that hold at every feasible point. Each step is recorded with its
source row, and replay repeats the sequence. The bounds are used only as
premises of certificates; the model keeps its declared bounds.

\subsection{Blocks, directions and activation}\label{sec:blocks}

Discovery runs at the first separator call, so a model solved during presolve
pays nothing. A ranged row $\ell\le g(v)\le u$ has two \emph{sides},
$g(v)\le u$ and $-g(v)\le-\ell$. Discovery writes each finite side of each
original row, including the objective epigraph row, exactly as a nonlinear
remainder plus an affine part. A side is admissible if its remainder involves
one to four variables with finite proven bounds and is either univariate or a
polynomial of total degree at most eight. Blocks are the variable sets of
single sides and greedy unions of overlapping sides, with at most four
variables, six nonlinear sides and sixteen signed affine domain sides; at
most 32 blocks are kept, in a fixed order with the largest first. Unsupported
remainders are declined as a whole, and a remainder is never altered. Affine
rows that involve only block variables, together with the proven bounds,
form the support domain $D$. Polynomial remainders of degree three or more in
three or four variables are admitted by discovery but are not certified by
any kernel (Bernstein bounds need one or two variables and at most 1{,}024
tensor coefficients), so such blocks yield no cuts.

At an LP solution the separator first tries, for each block and each of its
source sides $j$, the direction $a=0$, $\lambda=e_j$. This bounds only the
side's nonlinear remainder; the side's affine terms in the block variables
enter the cut only through the elimination. The support of the whole side,
with $a$ equal to the side's coefficients of the block variables, can be much
stronger; a variant that tries it first (\emph{whole-row directions}) is
evaluated in Section~\ref{sec:mechanism}. The
separator then solves up to three LPs that separate the current point, after
coordinate scaling, from the convex hull of sampled graph points plus the
cone of the side coordinates; an LP direction is used only if its scaled
violation exceeds $10^{-5}$. Samples are grid points of the box (33 points in
one variable, $7\times7$ in two, $3^d$ in $d=3,4$), plus the exact vertices of
the domain polytope and their centroid when the enumeration budget allows.
The binary64 values of $a$ and $\lambda$ are then certified. Quadratic blocks
are certified with Theorem~\ref{thm:polytope} when the number of row subsets
is at most 2{,}000; for four variables this allows at most seven signed
affine sides, and with at most sixteen sides the budget is never exceeded in
dimension three or less. Otherwise only the star oracle of
Theorem~\ref{thm:star} can apply, because the Bernstein bound is limited to
two variables and the ball-arithmetic bound to one; if the block is not a
star, that direction yields no cut. When an exact support computation returns
a minimizer that the samples missed, the minimizer is added and the LP is
solved again.

Two modes add cuts. Mode \code{all} tries the admitted blocks in their fixed
order until a cap is reached. Mode \code{auto} admits only quadratic blocks
that contain a side not proved convex and that have either an affine domain
row that is not a bound or at least two distinct nonlinear rows; it requires
a larger violation and stops trying a block after two failed certifications.
Both modes cap the number of callbacks, cuts and support calls
(Appendix~\ref{app:software}) and give all separator work a cumulative
allowance of $\min\{1\,\mathrm s,\,0.05\,T\}$ for a total budget $T$, checked
between operations; a single symbolic operation can overrun it. A cut is
added only if its violation divided by $\max\{1,\norm{\hat c}_1\}$ exceeds
$10^{-5}$ in \code{all} and $5\cdot10^{-4}$ in \code{auto}. Mode
\code{baseline} uses the same model builder and adds nothing. Modes with raised
limits and the variant with whole-row directions are listed in
Table~\ref{tab:modes}; we call the default first directions
\emph{remainder directions}.
